\documentclass[11pt,a4paper]{article}

\usepackage[utf8]{inputenc}
\usepackage[T1]{fontenc}
\usepackage{lmodern}
\usepackage[margin=2.5cm]{geometry}
\usepackage{parskip}
\usepackage{ragged2e}
\usepackage[hidelinks]{hyperref}

\pagestyle{empty}
\setlength{\parindent}{0pt}

\begin{document}

% Candidate information
{\Large\bfseries John Doe}\\[0.3cm]


\href{mailto:john@example.com}{john@example.com}\\



+33 6 00 00 00 00\\



Paris\\


\vspace{1cm}

% Recipient
\begin{flushright}

\textbf{Veolia}\\



Paris\\


\vspace{0.5cm}

12 novembre 2026

\end{flushright}

\vspace{0.8cm}

% Subject
\textbf{Subject: Candidature spontanée – Stage en Data Engineering}

\vspace{0.6cm}

% Salutation

Madame, Monsieur,


\vspace{0.5cm}

% Cover letter body
\justifying


Je me permets de vous adresser ma candidature spontanée pour un stage en Data Engineering au sein de Veolia, acteur du secteur de l’environnement à Paris. Étudiant en Master informatique, je souhaite mettre mes compétences techniques au service de vos équipes.

\par

Mes études et un stage en développement m’ont permis de pratiquer Python et SQL, ainsi que d’évoluer sous Linux. Je souhaite contribuer à des missions telles que l’écriture de scripts, l’interrogation de bases de données et l’automatisation, tout en consolidant mes bases en ingénierie des données au contact de vos experts.

\par

Je serais heureux d’échanger afin de préciser vos besoins et la manière dont je peux y contribuer.

\par


\vspace{0.7cm}

% Closing

Cordialement,


\vspace{0.5cm}

\textbf{John Doe}

\end{document}

